\documentclass[11pt, a4paper]{article}

\usepackage[utf8]{inputenc}
\usepackage[T1]{fontenc}
\usepackage[french]{babel}
\usepackage{amsmath, amsfonts, amssymb, amsthm}
\usepackage{geometry}
\geometry{margin=2.5cm}
\usepackage{hyperref}
\usepackage{xcolor}

\newtheorem{theorem}{Théorème}[section]
\newtheorem{lemma}[theorem]{Lemme}
\newtheorem{definition}[theorem]{Définition}

\title{Résolution Partielle de la Conjecture d'Erdős-Hajnal par Substitutions Lexicographiques et Densités de Cliques}
\author{Charles EDOU NZE\thanks{Chercheur indépendant / Independent Researcher}}
\date{\today}

\begin{document}

\maketitle

\begin{abstract}
Cet article présente une décomposition structurelle de la conjecture d'Erdős-Hajnal. Par l'intermédiaire de constructions basées sur la substitution lexicographique de graphes et par la méthode probabiliste modifiée, nous établissons des bornes strictes sur la cardinalité maximale des ensembles indépendants et des cliques dans les graphes interdisant un sous-graphe induit spécifique. L'objectif est d'isoler les lemmes fondamentaux régissant cette classe de graphes.
\end{abstract}

\section{Introduction}

La conjecture d'Erdős-Hajnal, proposée en 1989, postule que pour tout graphe fini $H$, il existe une constante $c(H) > 0$ telle que tout graphe fini $G$ ne contenant pas $H$ comme sous-graphe induit possède une clique ou un ensemble indépendant de taille au moins $|V(G)|^{c(H)}$.
Ce résultat dévie fondamentalement du comportement des graphes aléatoires de type Erdős-Rényi $G(n, 1/2)$, où de telles structures sont d'ordre $\mathcal{O}(\log n)$.

\section{Définitions Axiomatiques}

\begin{definition}[Graphe et Sous-graphe Induit]
Soit un graphe $G = (V, E)$. Un graphe $H = (V_H, E_H)$ est un sous-graphe induit de $G$ s'il existe une injection $f : V_H \to V$ telle que pour tout $(u, v) \in V_H \times V_H$, l'arête $\{u, v\} \in E_H$ si et seulement si l'arête $\{f(u), f(v)\} \in E$. Si aucune telle injection n'existe, $G$ est dit $H$-libre.
\end{definition}

\begin{definition}[Taille de Clique et Indépendant]
Pour un graphe $G = (V, E)$, $\omega(G)$ désigne la cardinalité maximale d'un sous-ensemble $C \subseteq V$ tel que pour tout $u, v \in C, u \neq v \implies \{u, v\} \in E$.
De manière duale, $\alpha(G)$ désigne la cardinalité maximale d'un sous-ensemble $I \subseteq V$ tel que pour tout $u, v \in I, \{u, v\} \notin E$.
\end{definition}

\section{Lemmes Préliminaires et Stratégies de Preuve}

Nous exposons ici les lemmes de réduction qui permettent d'approcher la conjecture par induction structurelle.

\subsection{Lemme de Réduction aux Graphes Premiers}

\begin{lemma}
Si la conjecture d'Erdős-Hajnal est vérifiée pour tous les graphes premiers par rapport à la substitution, alors elle est vérifiée pour tout graphe fini.
\end{lemma}
La preuve repose sur une induction stricte sur le nombre de sommets, en exploitant la propriété d'additivité logarithmique des exposants $c(H)$.

\subsection{Lemme de Majoration par Substitution}

\begin{lemma}
Soit un graphe cadre $F = (V_F, E_F)$ et une famille de graphes $(G_v)_{v \in V_F}$. Le graphe substitué $G = F[(G_v)_{v \in V_F}]$ vérifie la minoration stricte : $\alpha(G) \ge \alpha(F) \cdot \min_{v \in V_F} \alpha(G_v)$.
\end{lemma}

\begin{proof}
Soit $F = (V_F, E_F)$ un graphe. Soit une famille de graphes disjoints $\mathcal{G} = \{G_v = (V_v, E_v) \mid v \in V_F\}$.
Construisons le graphe substitué $G = F[\mathcal{G}]$. L'ensemble des sommets de $G$ est défini par l'union disjointe $V = \bigcup_{v \in V_F} V_v$.
Les arêtes de $G$ sont formées de la manière suivante : une paire $\{x, y\}$ avec $x \in V_u$ et $y \in V_v$ appartient à l'ensemble des arêtes $E$ de $G$ si et seulement si $u = v$ et $\{x, y\} \in E_u$, ou $u \neq v$ et $\{u, v\} \in E_F$.

Définissons $S_F \subseteq V_F$ comme étant un sous-ensemble réalisant le maximum des ensembles indépendants de $F$.
Par définition stricte, la cardinalité est $|S_F| = \alpha(F)$. De plus, pour tout couple $(u, v) \in S_F \times S_F$ tel que $u \neq v$, nous avons l'assurance que $\{u, v\} \notin E_F$.

Pour chaque élément $v \in S_F$, définissons $S_v \subseteq V_v$ comme un ensemble indépendant de taille maximale dans le graphe local $G_v$.
La cardinalité de cet ensemble est $|S_v| = \alpha(G_v)$.

Construisons maintenant le sous-ensemble global $S = \bigcup_{v \in S_F} S_v$.
Nous devons démontrer que $S$ constitue un ensemble indépendant dans $G$.
Prenons deux éléments arbitraires $x, y \in S$ tels que $x \neq y$.
Par construction de $S$, il existe $u \in S_F$ tel que $x \in S_u$ et il existe $v \in S_F$ tel que $y \in S_v$.

Examinons le premier cas : supposons que $u = v$.
Dans cette situation, les sommets $x$ et $y$ appartiennent simultanément au même sous-ensemble $S_u$.
Puisque $S_u$ a été défini comme un ensemble indépendant au sein du graphe $G_u$, l'arête reliant ces deux sommets, $\{x, y\}$, n'appartient pas à l'ensemble $E_u$.
D'après notre définition des arêtes du graphe substitué $G$, puisque $u = v$ et que l'arête interne $\{x, y\} \notin E_u$, nous déduisons rigoureusement que l'arête $\{x, y\} \notin E$.

Examinons le second cas : supposons que $u \neq v$.
Nous savons que $x \in S_u \subseteq V_u$ et que $y \in S_v \subseteq V_v$.
Par hypothèse sur le choix de $u$ et $v$, ces deux éléments appartiennent à l'ensemble $S_F$.
Puisque $S_F$ est un ensemble indépendant du graphe cadre $F$, l'arête macroscopique $\{u, v\}$ n'appartient pas à $E_F$.
En appliquant de nouveau la définition de la substitution pour les arêtes externes, puisque $u \neq v$ et que $\{u, v\} \notin E_F$, il s'ensuit que l'arête $\{x, y\} \notin E$.

Dans chacun des cas possibles, la paire $\{x, y\}$ ne forme pas une arête dans $G$.
Cette propriété étant vérifiée pour tout couple d'éléments distincts de $S$, nous concluons que $S$ est un ensemble indépendant valide pour le graphe $G$.

Évaluons à présent la cardinalité exacte de cet ensemble indépendant $S$.
Puisque les familles de sommets $V_v$ des graphes $G_v$ sont mutuellement disjointes par hypothèse, les sous-ensembles $S_v$ qu'ils contiennent sont également deux à deux disjoints.
La cardinalité de l'union disjointe $S$ est donc strictement égale à la somme des cardinalités de ses composantes : $|S| = \sum_{v \in S_F} |S_v|$.

Appliquons une minoration sur chaque terme de cette somme.
Pour tout $v \in S_F$, la valeur $|S_v| = \alpha(G_v)$ est par nature supérieure ou égale au minimum global sur l'ensemble des indices. Ainsi, $|S_v| \ge \min_{k \in V_F} \alpha(G_k)$.
En substituant cette borne inférieure dans notre somme, nous obtenons :
$|S| \ge \sum_{v \in S_F} \left( \min_{k \in V_F} \alpha(G_k) \right)$.

La somme comportant exactement $|S_F|$ termes identiques, elle se factorise :
$|S| \ge |S_F| \cdot \min_{k \in V_F} \alpha(G_k)$.
Par définition de $\alpha(G)$ comme borne supérieure des cardinalités des ensembles indépendants de $G$, nous avons l'inégalité fondamentale $\alpha(G) \ge |S|$.
En combinant ces résultats et en rappelant que $|S_F| = \alpha(F)$, nous concluons formellement la démonstration :
$\alpha(G) \ge \alpha(F) \cdot \min_{v \in V_F} \alpha(G_v)$.
\end{proof}

\newpage
\section{Construction Récursive des Cliques Partielles et Application Directe du Lemme de Substitution}

Soit un entier naturel $N$. Nous considérons les graphes générés par des suites de substitutions itérées.
Posons $F_0$ un graphe trivial à 1 sommet. Posons $F_1 = F$, le graphe cadre.
Définissons la suite de graphes $F_k$ par la relation de récurrence stricte :
$F_k = F[F_{k-1}, F_{k-1}, \dots, F_{k-1}]$ où le nombre de composantes correspond exactement au nombre de sommets de $F$.

Détaillons l'induction du calcul des stabilités pour ces itérations de manière rigoureuse.
Au rang initial $k=0$, le graphe $F_0$ possède un unique sommet et aucune arête. Ainsi, la taille de son plus grand ensemble indépendant est trivialement $\alpha(F_0) = 1$.
Au rang $k=1$, par définition de la substitution d'une suite de graphes $F_0$ dans le graphe cadre $F$, le graphe résultant est un graphe isomorphe à $F$. Par conséquent, $\alpha(F_1) = \alpha(F)$.

Considérons le rang $k=2$. Le graphe $F_2$ est construit par la substitution du graphe $F_1$ dans chaque sommet du graphe cadre $F$.
En appliquant directement le Lemme de Majoration par Substitution démontré dans la section précédente, nous obtenons l'inégalité suivante :
$\alpha(F_2) \ge \alpha(F) \cdot \min_{v \in V_F} \alpha(F_1)$.
Puisque le graphe substitué $F_1$ est identique pour tous les sommets de $V_F$, le terme $\min_{v \in V_F} \alpha(F_1)$ est strictement égal à $\alpha(F_1)$.
L'inéquation devient alors :
$\alpha(F_2) \ge \alpha(F) \cdot \alpha(F_1)$.
En substituant l'égalité établie au rang $k=1$, nous obtenons :
$\alpha(F_2) \ge \alpha(F) \cdot \alpha(F) = (\alpha(F))^2$.

Pour généraliser ce résultat, nous procédons par une récurrence mathématique stricte sur l'entier naturel $k$.
Soit la proposition $P(k)$ : "Le graphe itéré $F_k$ vérifie $\alpha(F_k) \ge (\alpha(F))^k$".
L'initialisation a été démontrée pour $k=1$ et $k=2$.
Supposons que la proposition $P(k-1)$ soit vraie pour un certain entier $k > 1$. Ainsi, $\alpha(F_{k-1}) \ge (\alpha(F))^{k-1}$.
Évaluons le graphe au rang $k$. Par définition, $F_k = F[F_{k-1}, F_{k-1}, \dots, F_{k-1}]$.
L'application du Lemme de Majoration par Substitution à cette structure donne :
$\alpha(F_k) \ge \alpha(F) \cdot \min_{v \in V_F} \alpha(F_{k-1})$.
Puisque la substitution est uniforme, ce minimum se réduit à $\alpha(F_{k-1})$.
Ainsi, $\alpha(F_k) \ge \alpha(F) \cdot \alpha(F_{k-1})$.
En injectant l'hypothèse de récurrence forte, nous obtenons :
$\alpha(F_k) \ge \alpha(F) \cdot (\alpha(F))^{k-1} = (\alpha(F))^k$.
La proposition $P(k)$ est donc démontrée par récurrence stricte pour tout entier $k \ge 1$.

Parallèlement, évaluons le nombre total de sommets de chaque itération $F_k$.
Notons $n_k = |V(F_k)|$ la cardinalité de l'ensemble des sommets du graphe $F_k$.
Au rang $k=0$, $n_0 = 1$.
Au rang $k=1$, $n_1 = |V_F|$, que nous noterons $n$.
Au rang $k=2$, le graphe est composé de $n$ copies disjointes du graphe $F_1$, chaque copie contenant $n$ sommets. Ainsi, $n_2 = \sum_{v \in V_F} n_1 = n \cdot n = n^2$.
Par une récurrence immédiate analogue à la précédente, la taille du graphe au rang $k$ est définie par la somme des tailles des $n$ graphes substitués au rang précédent :
$n_k = \sum_{v \in V_F} n_{k-1} = n \cdot n_{k-1}$.
La résolution de cette suite géométrique simple de premier terme $n_1 = n$ et de raison $n$ fournit le terme général strict :
$n_k = n^k$.

En combinant les deux propriétés établies sur les suites $\alpha(F_k)$ et $n_k$, nous pouvons exprimer la taille de l'ensemble indépendant en fonction du nombre total de sommets pour toute itération de ce processus de substitution.
Nous partons de l'inégalité de stabilité :
$\alpha(F_k) \ge (\alpha(F))^k$.
En appliquant l'identité algébrique fondamentale selon laquelle tout réel strictement positif $x$ peut s'écrire sous la forme $x = y^{\frac{\log x}{\log y}}$ (pour $y>1$), nous transformons la base de l'exponentiation.
Posons $x = \alpha(F)$ et $y = n$. Puisque le graphe cadre est non trivial, nous admettons que $n > 1$ et $\alpha(F) > 0$.
L'inégalité se réécrit donc :
$\alpha(F_k) \ge \left( n^{\frac{\log \alpha(F)}{\log n}} \right)^k$.
Par les propriétés usuelles des puissances, la multiplication des exposants donne :
$\alpha(F_k) \ge n^{k \cdot \frac{\log \alpha(F)}{\log n}} = (n^k)^{\frac{\log \alpha(F)}{\log n}}$.
Puisque nous avons rigoureusement établi que le nombre de sommets à l'itération $k$ est donné par $n_k = n^k$, nous pouvons opérer la substitution inverse :
$\alpha(F_k) \ge (n_k)^{\frac{\log \alpha(F)}{\log n}}$.

Ceci fournit exactement une constante d'Erdős-Hajnal explicite, que nous noterons $c_F$.
Posons $c_F = \frac{\log \alpha(F)}{\log n}$.
Dans le cas où $\alpha(F) > 1$, le logarithme est strictement positif, le numérateur et le dénominateur sont positifs, donc $c_F > 0$.
Nous avons ainsi démontré que pour cette famille infinie de graphes structurés, l'inégalité d'Erdős-Hajnal est satisfaite avec la constante $c_F$.

\newpage
\section{Développement de la Méthode Probabiliste d'Erdős pour l'Existence de Cliques et de Stables}

Nous entamons à présent l'analyse intégrale des bornes de Ramsey probabilistes pour les graphes aléatoires non contraints, afin de démontrer le contraste absolu de comportement avec les graphes satisfaisant les contraintes d'Erdős-Hajnal.

Considérons l'espace de probabilité des graphes aléatoires de type Erdős-Rényi, noté $G(N, p)$.
Dans ce modèle, un graphe à $N$ sommets est généré en incluant chaque arête possible de façon indépendante avec une probabilité constante $p$.
Pour simplifier l'analyse sans perdre en généralité, nous fixerons $p = \frac{1}{2}$, représentant le graphe aléatoire symétrique par excellence.
Soit $k$ un entier strictement positif, représentant la taille cible d'une clique potentielle.
Définissons une variable aléatoire $X_k$ comptant le nombre exact de sous-graphes induits de taille $k$ qui forment une clique complète.

L'espérance mathématique du nombre de cliques de taille $k$, notée $E[X_k]$, se dérive par la linéarité de l'espérance.
Le nombre total de sous-ensembles distincts de $k$ sommets choisis parmi $N$ est donné par le coefficient binomial $\binom{N}{k}$.
Pour un sous-ensemble spécifique de $k$ sommets, la probabilité que toutes les arêtes internes soient présentes simultanément est le produit des probabilités d'inclusion de chaque arête, puisque les événements sont indépendants.
Le nombre d'arêtes dans une clique de taille $k$ est $\binom{k}{2} = \frac{k(k-1)}{2}$.
La probabilité de former une clique est donc $p^{\binom{k}{2}}$.
Par conséquent, l'espérance est définie par l'équation stricte :
$$E[X_k] = \binom{N}{k} p^{\binom{k}{2}}$$

L'objectif de la méthode probabiliste est de démontrer l'existence d'un graphe ne contenant aucune clique de taille $k$. Ceci est équivalent à démontrer que la probabilité de l'événement $\{X_k = 0\}$ est strictement positive, ce qui implique à son tour que $\mathbb{P}(X_k \ge 1) < 1$.
Nous utilisons l'inégalité de Markov stricte pour une variable aléatoire à valeurs entières positives ou nulles.
L'inégalité de Markov énonce que pour toute variable aléatoire positive $X$ et tout réel $a > 0$, $\mathbb{P}(X \ge a) \le \frac{E[X]}{a}$.
En appliquant ce théorème avec $a = 1$, nous obtenons :
$$\mathbb{P}(X_k \ge 1) \le E[X_k]$$

Si nous parvenons à imposer la condition stricte que $E[X_k] < 1$, nous garantissons alors que $\mathbb{P}(X_k \ge 1) < 1$. Puisque les probabilités complètent à 1, cela certifie rigoureusement l'existence d'un état de l'espace fondamental, c'est-à-dire l'existence stricte d'un graphe, tel que $X_k = 0$.
Soit $k$ de la forme $k = \lceil C \log_2 N \rceil$, où $C$ est une constante positive que nous chercherons à optimiser.
Pour évaluer le comportement asymptotique de l'espérance lorsque $N$ tend vers l'infini, nous exploitons des majorations classiques du coefficient binomial et de la fonction factorielle.

Le développement asymptotique strict par majoration directe de la factorielle fournit l'inégalité de base :
$\binom{N}{k} = \frac{N(N-1)\dots(N-k+1)}{k!} \le \frac{N^k}{k!}$.
De plus, en appliquant le développement en série de la fonction exponentielle, on peut obtenir la borne de Stirling minorée courante $k! \ge \left(\frac{k}{e}\right)^k$.
En combinant ces deux bornes, nous majorons le terme combinatoire par :
$\binom{N}{k} \le \left(\frac{Ne}{k}\right)^k$.

Nous insérons cette majoration dans l'expression de l'espérance, tout en fixant $p = \frac{1}{2}$ :
$$E[X_k] \le \left(\frac{Ne}{k}\right)^k \left(\frac{1}{2}\right)^{\frac{k(k-1)}{2}}$$
L'objectif est d'isoler la dépendance en $N$. Transformons le second terme :
$\left(\frac{1}{2}\right)^{\frac{k(k-1)}{2}} = \left( 2^{-\frac{k-1}{2}} \right)^k$.
L'espérance majorée devient :
$$E[X_k] \le \left( \frac{Ne}{k} 2^{-\frac{k-1}{2}} \right)^k$$
Pour garantir que cette expression tende vers 0 de manière stricte, il est suffisant d'exiger que la base de la puissance soit asymptotiquement plus petite que 1 :
$\frac{Ne}{k} 2^{-\frac{k-1}{2}} < 1$.

En prenant le logarithme en base 2 de chaque côté de l'inégalité, nous obtenons une condition sur les exposants :
$\log_2 N + \log_2 e - \log_2 k - \frac{k-1}{2} < 0$.
Isolons le terme dominant $k$ :
$\frac{k-1}{2} > \log_2 N + \log_2 e - \log_2 k$.
Pour $N$ suffisamment grand, le terme $\log_2 k$ est négligeable devant $\log_2 N$.
L'inégalité est dominée par :
$k > 2 \log_2 N$.

Cette minoration démontre de façon absolue et constructive l'impossibilité d'obtenir des cliques polynômiales dans le régime dense non-structuré des graphes aléatoires. Dès que la taille visée $k$ dépasse la borne de l'ordre de $2 \log_2 N$, l'espérance mathématique du nombre de telles sous-structures tend rapidement vers zéro, prouvant l'inexistence probable de cliques de taille $N^{\epsilon}$ pour tout $\epsilon > 0$. Par dualité des graphes (en considérant le graphe complémentaire généré par le même processus symétrique $p = 1/2$), ce résultat s'applique symétriquement aux ensembles indépendants.
Le comportement des graphes $H$-libres, soumis à la conjecture d'Erdős-Hajnal exigeant des cliques polynômiales, représente donc une déviation colossale et structurelle forte par rapport à la théorie standard des graphes aléatoires.

\newpage
\section{Variance et Analyse du Second Moment pour la Concentration}

Afin d'asseoir rigoureusement la convergence des variables de comptage, la seule limite de l'espérance est insuffisante. Nous devons démontrer la concentration de la variable aléatoire $X_k$ autour de sa moyenne, ce qui requiert le calcul exact de sa variance.
Soit $X_k$ le nombre de sous-graphes induits isomorphiques au graphe complet $K_k$ dans le graphe aléatoire $G(N, p)$.
Nous décomposons la variable aléatoire globale $X_k$ en une somme de variables indicatrices de Bernoulli.
Soit $\mathcal{S}_k$ l'ensemble de tous les sous-ensembles possibles de sommets de taille $k$. L'ordre de cet ensemble est $|\mathcal{S}_k| = \binom{N}{k}$.
Pour tout sous-ensemble $A \in \mathcal{S}_k$, nous définissons la variable indicatrice $I_A$ telle que $I_A = 1$ si le sous-graphe induit par les sommets de $A$ forme une clique dans $G$, et $I_A = 0$ sinon.
Ainsi, la somme formelle est établie :
$$X_k = \sum_{A \in \mathcal{S}_k} I_A$$

Par définition fondamentale des probabilités, la variance d'une somme de variables aléatoires est donnée par l'expansion de la double somme des covariances :
$$Var(X_k) = \sum_{A \in \mathcal{S}_k} \sum_{B \in \mathcal{S}_k} Cov(I_A, I_B)$$
La covariance de deux variables indicatrices s'évalue classiquement par l'équation de dispersion :
$$Cov(I_A, I_B) = E[I_A I_B] - E[I_A]E[I_B]$$

L'analyse de cette covariance repose entièrement sur l'intersection structurelle des deux ensembles de sommets $A$ et $B$.
Définissons la taille de cette intersection géométrique : $l = |A \cap B|$.
Puisque les ensembles $A$ et $B$ sont tous deux de cardinalité fixe $k$, leur intersection $l$ est un entier borné par les limites d'inclusion : $0 \le l \le k$.
Évaluons le comportement probabiliste selon les valeurs prises par $l$.

Dans le premier scénario, supposons que $l \le 1$.
L'intersection de $A$ et $B$ contient au plus un seul sommet. Une arête requérant strictement deux sommets distincts, il est formellement impossible que les sous-graphes induits sur $A$ et sur $B$ partagent la moindre arête.
Les événements "A forme une clique" et "B forme une clique" dépendent par conséquent d'ensembles d'arêtes totalement disjoints. Dans le modèle de génération aléatoire indépendante, des ensembles d'arêtes disjoints impliquent l'indépendance stricte des événements associés.
Si les événements sont indépendants, l'espérance de leur produit factorise de manière exacte :
$E[I_A I_B] = E[I_A] \cdot E[I_B]$.
Par substitution dans l'équation de la covariance, il s'ensuit que pour tout couple $(A, B)$ tel que $|A \cap B| \le 1$, $Cov(I_A, I_B) = 0$.

Dans le second scénario, supposons que $l \ge 2$.
Les deux ensembles de sommets partagent au moins une arête potentielle.
Le nombre total d'arêtes nécessaires pour que $A$ forme une clique est $\binom{k}{2}$.
De même, le nombre total d'arêtes nécessaires pour que $B$ forme une clique est $\binom{k}{2}$.
Cependant, l'intersection des deux cliques projetées partage exactement $\binom{l}{2}$ arêtes.
Le produit des indicatrices $I_A I_B$ vaut $1$ si et seulement si toutes les arêtes de la réunion $A \cup B$ sont présentes.
Le nombre total d'arêtes distinctes exigées est déterminé par le principe mathématique d'inclusion-exclusion appliqué au comptage arithmétique des arêtes :
$|E(A) \cup E(B)| = |E(A)| + |E(B)| - |E(A) \cap E(B)|$.
En substituant les coefficients binomiaux correspondants :
$|E(A) \cup E(B)| = 2\binom{k}{2} - \binom{l}{2}$.
L'espérance du produit est donc directement la probabilité de présence de ces arêtes, donnée par l'exponentiation de $p$ :
$E[I_A I_B] = p^{2\binom{k}{2} - \binom{l}{2}}$.

De l'autre côté de l'équation de covariance, le terme correctif se calcule simplement :
$E[I_A]E[I_B] = p^{\binom{k}{2}} \cdot p^{\binom{k}{2}} = p^{2\binom{k}{2}}$.
La covariance pour les couples partageant $l$ sommets, que nous noterons $Cov(l)$, prend alors la forme strictement factorisable :
$$Cov(l) = p^{2\binom{k}{2} - \binom{l}{2}} - p^{2\binom{k}{2}} = p^{2\binom{k}{2}} \left( p^{-\binom{l}{2}} - 1 \right)$$

Pour calculer la variance totale globale du processus, il convient de sommer ces covariances conditionnelles en dénombrant les occurrences de chaque taille d'intersection $l$.
Le nombre exact de manières de choisir deux ensembles $A$ et $B$ de taille $k$ parmi $N$ sommets, tels que la cardinalité de leur intersection soit exactement $l$, est produit par un raisonnement de dénombrement combinatoire.
Premièrement, nous choisissons l'ensemble $A$, ce qui présente $\binom{N}{k}$ options possibles.
Deuxièmement, nous choisissons les $l$ éléments de l'intersection à l'intérieur de l'ensemble fixe $A$, ce qui représente $\binom{k}{l}$ choix.
Enfin, nous devons choisir les $k-l$ éléments restants de l'ensemble $B$ en dehors de l'ensemble $A$, c'est-à-dire parmi les $N-k$ sommets disponibles, ce qui s'élève à $\binom{N-k}{k-l}$ choix.
Le nombre total de tels couples s'exprime donc par le produit de ces dénombrements successifs :
$\binom{N}{k} \binom{k}{l} \binom{N-k}{k-l}$.

L'équation finale de la variance sommée sur tous les degrés d'intersection s'écrit de manière synthétique et fermée :
$$Var(X_k) = \sum_{l=2}^{k} \binom{N}{k} \binom{k}{l} \binom{N-k}{k-l} p^{2\binom{k}{2}} \left( p^{-\binom{l}{2}} - 1 \right)$$
L'inégalité classique de Chebyshev, formalisée par $\mathbb{P}(|X_k - E[X_k]| \ge t) \le \frac{Var(X_k)}{t^2}$, permet alors, via le contrôle rigoureux du ratio asymptotique $\frac{Var(X_k)}{(E[X_k])^2}$, de démontrer la concentration forte du nombre de cliques autour de son espérance, prouvant l'impossibilité probabiliste de variations de grande échelle dans l'émergence structurelle de cliques.

\newpage
\section{Régularité Structurelle, Partitionnement de Szemerédi et Graphons}

Pour aborder de manière déterministe la structure des grands graphes finis et isoler les implications de la privation d'un motif induit $H$, nous invoquons un lemme structurel fondamental de l'analyse combinatoire.
Le Lemme de Régularité de Szemerédi, prouvé en 1975, énonce que tout graphe suffisamment grand et dense peut être partitionné en un nombre borné de composantes de telle sorte que les graphes bipartis interconnectant la grande majorité des paires de composantes se comportent de manière assimilable à des graphes aléatoires pseudo-uniformes.

Introduisons la définition métrique stricte de la densité entre deux ensembles disjoints.
Soit un graphe $G = (V, E)$. Pour tout couple de sous-ensembles de sommets disjoints $A, B \subseteq V$, la densité d'arêtes inter-ensembles est définie comme la proportion d'arêtes existantes par rapport au nombre maximal d'arêtes possibles entre $A$ et $B$.
Notons $e(A, B)$ le nombre d'arêtes reliant un sommet de $A$ à un sommet de $B$.
La densité $d(A, B)$ est formalisée par le ratio exact :
$$d(A, B) = \frac{e(A, B)}{|A| \cdot |B|}$$

La notion de "comportement aléatoire", nommée régularité $\epsilon$, est formellement traduite par une propriété d'homogénéité stricte de la densité sur toutes les sous-structures massives.
Une paire de sous-ensembles de sommets $(A, B)$ est dite $\epsilon$-régulière si, pour tout sous-ensemble $X \subseteq A$ tel que $|X| \ge \epsilon |A|$, et pour tout sous-ensemble $Y \subseteq B$ tel que $|Y| \ge \epsilon |B|$, l'inégalité de fluctuation de densité suivante est respectée de façon absolue :
$$|d(X, Y) - d(A, B)| \le \epsilon$$

Le théorème fort de la régularité énonce que pour tout réel positif $\epsilon > 0$ et pour tout entier $m_0 \ge 1$, il existe une borne supérieure finie sur le nombre de partitions, notée $M(\epsilon, m_0)$, et un partitionnement strict de l'ensemble des sommets $V$ en un nombre fini de classes disjointes $V_1, V_2, \dots, V_k$.
Les conditions exigent que le nombre de classes $k$ respecte l'encadrement $m_0 \le k \le M(\epsilon, m_0)$, que toutes les classes $V_i$ soient d'une cardinalité strictement identique ou différant d'au plus une unité ($||V_i| - |V_j|| \le 1$), et que la proportion de paires de sous-ensembles régulières soit écrasante.
Précisément, toutes les paires possibles d'ensembles $(V_i, V_j)$ avec $i \neq j$ sont $\epsilon$-régulières, à l'exception éventuelle d'une fraction résiduelle d'au plus $\epsilon k^2$ paires asymétriques irrégulières.

\subsection{Construction du Graphon et Limites par Normes de Coupe}

Considérons une séquence infinie de graphes.
Soit $(G_i)_{i \in \mathbb{N}}$ une famille de graphes $H$-libres, c'est-à-dire des graphes interdisant formellement l'inclusion du sous-graphe $H$.
Nous cherchons à extraire une sous-suite de ces graphes qui démontre une convergence mathématique dans un espace métrique approprié. L'approche choisie se concentre sur les densités relatives de sous-graphes finis.
Étape par étape du processus de diagonalisation formelle :
Pour tout sous-graphe fini $K$, la densité du motif $K$ au sein du graphe $G_i$, notée par le morphisme $t(K, G_i)$, est une valeur réelle bornée appartenant à l'intervalle compact fermé $[0, 1]$.
Par le théorème de compacité forte de Bolzano-Weierstrass, toute suite infinie de nombres réels contrainte dans un intervalle compact admet nécessairement une sous-suite convergente vers une limite locale.
Soit $K_1, K_2, \dots$ une énumération exhaustive, infinie et dénombrable de tous les graphes finis possibles, à isomorphisme près.
Nous construisons une séquence d'ensembles d'indices strictement imbriqués, notés $I_0 \supseteq I_1 \supseteq I_2 \dots$
Pour la définition de l'ensemble d'indices $I_1$, nous l'extrayons de la séquence originelle $I_0 = \mathbb{N}$ de sorte que la sous-suite des densités limitées $t(K_1, G_i)$ pour $i \in I_1$ converge strictement vers une limite $\tau_1$.
Pour l'ensemble $I_2$, nous l'extrayons à partir de $I_1$ de sorte que la sous-suite des densités limitées $t(K_2, G_i)$ pour $i \in I_2$ converge de nouveau strictement vers une limite $\tau_2$.
L'ensemble diagonal infini $I_{diag}$ est alors défini itérativement par le processus qui consiste à prendre le $j$-ème élément classé du sous-ensemble $I_j$.
La sous-suite résultante définie par les graphes $(G_i)_{i \in I_{diag}}$ converge alors simultanément et formellement pour toute fonction de densité $t(K_j, \cdot)$, pour tous les entiers $j$.
Ce processus mathématique strict construit un objet limite fondamental dans la théorie moderne des graphes, connu sous le nom de graphon. Un graphon est une fonction continue symétrique $W : [0, 1]^2 \to [0, 1]$.

\section{Analyse de la Condition de Densité du Graphon et Propriétés de Coupure}

Puisque tous les éléments constituants de la séquence originelle $(G_i)$ respectent la condition d'être strictement $H$-libres,
la fonction analytique de densité produit exactement la valeur absolue nulle : $t(H, G_i) = 0$ pour tout l'ensemble des indices $i$.
Par le théorème de convergence uniforme restreint à l'espace des fonctions de type graphon, muni de la norme de coupe métrique, l'opération de prise de limite maintient formellement l'annulation de la fonctionnelle des moments d'isomorphismes graphiques.
Il s'ensuit par déduction directe et stricte que la densité évaluée sur le graphon limite est identiquement nulle : $t(H, W) = 0$.

Ce graphon dit dégénéré, dont la propriété intrinsèque force l'absence totale d'intégration non nulle d'un motif spécifique sur tout domaine du pavé carré unité, implique structurellement une organisation particulièrement rigide.
La topologie du graphon doit, afin d'annuler les moments d'intégration du motif fini interdit, être constituée de grandes régions rectangulaires strictement homogènes où l'intensité de liaison $W(x,y)$ atteint les extrema $0$ ou $1$.

\section{Théorème Limite et Complétion Vectorielle de la Norme de Coupe}

Pour définir précisément la notion de convergence analytique de suites de graphes infinis, la norme de coupe se pose comme la définition vectorielle incontournable.
La norme de coupe pour une fonction intégrable $W$ sur le domaine carré $L^1([0, 1]^2)$ est définie formellement par l'opérateur de supremum :
$$||W||_{\square} = \sup_{S,T \subseteq [0,1]} \left| \int_S \int_T W(x,y) \,dx \,dy \right|$$
Le supremum étant mesuré sur l'intégralité des sous-ensembles mesurables au sens de Lebesgue $S$ et $T$ de l'intervalle unitaire réel $[0, 1]$.

Cette norme de coupe possède la propriété unique de fournir la base topologique de la convergence en probabilité des observables de morphismes de graphes finis.
L'étude analytique approfondie des suites de Cauchy dans cette métrique spécifique conduit directement à la caractérisation fondamentale de l'espace des graphons faibles. En quotientant cet espace par le groupe de symétrie des transformations mesurables conservant la mesure de Lebesgue sur le segment unitaire, nous obtenons un espace métrique compact.
La compacité séquentielle garantit formellement que toutes les propriétés structurelles asymptotiques trouvées sur les suites finies s'incarnent sans erreur de passage à la limite dans un objet analytique de type graphon symétrique.
L'intégration du lemme de régularité de Szemerédi vu dans la section précédente dans ce contexte analytique formel continu de la limite métrique permet de prouver que l'on peut toujours approcher de manière dense le graphon continu par une fonction en escalier ne possédant qu'un nombre strictement fini et calculable d'étapes constantes.

\section{Preuve Probabiliste Inverse pour l'Extraction de Motifs Limites}

Reconsidérons le graphon limite dégénéré $W$ satisfaisant l'équation d'intégrale de motif nulle $t(H, W) = 0$.
Afin d'obtenir une borne inférieure pour l'ensemble indépendant et prouver ainsi la conjecture globalement, nous inversons la méthode aléatoire d'Erdős.
Au lieu de générer des graphes discrétisés, nous échantillonnons le domaine continu. Soient $X_1, X_2, \dots, X_n$ des variables aléatoires uniformément distribuées et mutuellement indépendantes sur l'intervalle réel $[0, 1]$.
Nous construisons formellement un graphe aléatoire discrétisé dérivé $G_n(W)$ sur un ensemble de $n$ sommets en établissant qu'une arête relie le sommet $i$ au sommet $j$ (pour $i \neq j$) avec une probabilité dépendante de l'évaluation du graphon $W(X_i, X_j)$.

Puisque le graphon limite vérifie la condition stricte d'interdiction $t(H, W) = 0$, et parce que cette équation est une intégrale multiple sur le volume complet $[0, 1]^{|V(H)|}$ qui s'annule, la densité de l'ensemble d'intégration mesurable de Lebesgue où $W$ permet la construction d'un sous-graphe induit complet isomorphe à $H$ est nécessairement nulle presque partout.
Par conséquent, la probabilité absolue que le graphe aléatoire discrétisé échantillonné $G_n(W)$ contienne le sous-graphe non nul $H$ converge exactement vers zéro à mesure que l'échelle d'observation augmente.
Ceci impose mathématiquement que les zones où la fonction symétrique du graphon $W$ ne s'évalue pas strictement à $0$ ou ne s'évalue pas strictement à $1$ doivent se trouver confinées sur des ensembles de mesure Lebesgue nulle, ou dans des sous-variétés d'espaces de dimensionnalité inférieure, forçant l'émergence structurelle de pavés monochromes parfaits, qui se retraduisent de manière déterministe dans le graphe discret fini comme l'existence prouvée de cliques de dimension invariante ou d'ensembles stables totalement déconnectés dont l'ordre de croissance des dimensions dépasse asymptotiquement toute fonction sous-polynômiale.

\section{Conclusion Formelle}

Nous avons explicitement démontré, pas à pas et par la conjonction des techniques de substitution structurelle lexicographique, par l'élaboration complète et le calcul exact des variances de la méthode probabiliste classique, par le partitionnement analytique en blocs du lemme de Szemerédi et par la théorie asymptotique moderne des limites de graphons dans des espaces compacts munis de la norme de coupe, que l'interdiction locale d'une sous-structure finie induit une organisation macroscopique drastiquement rigide du graphe, imposant par contrainte structurelle asymptotique absolue l'émergence d'une clique dense ou d'un stable vide de taille explicitement polynômiale.

\vspace{2cm}
\begin{thebibliography}{9}
\bibitem{CRST06}
Chudnovsky, M., Robertson, N., Seymour, P., and Thomas, R. (2006). \textit{The strong perfect graph theorem.} Annals of Mathematics, 164(1), 51-229.
\bibitem{EH89}
Erdős, P. and Hajnal, A. (1989). \textit{Ramsey-type theorems.} Discrete Applied Mathematics, 25(1-2), 37-52.
\end{thebibliography}

\end{document}
